\pagenumbering{arabic}

\section{Einleitung}
\label{sec:einleitung}

%% TODO: Sobald die Aufteilung der Kapitel im Team feststeht, hier ggf.
%% \subsection{... -- Name} je Verfasser:in ergänzen (siehe Vorlage),
%% analog zur Vorgabe aus der Aufgabenstellung, individuelle Beiträge
%% zu kennzeichnen.

Die Trainingszelle \textit{KUKA ready2\_educate\_pro} des Fachbereichs 2 wird im
Rahmen der Lehrveranstaltung \enquote{Simulation und Regelung} für praxisnahe
Projekte im Studiengang Mechatronik und Robotik eingesetzt. Zwei vorangegangene
Abschluss- bzw.\ Projektarbeiten haben die Zelle bereits für eigene Zwecke
angebunden. Friesen~\cite{friesen2021} realisierte eine Verbindung zwischen
MATLAB und der KR~C4-Steuerung über den \ac{kvp}-Dienst, um Bahnen aus MATLAB
heraus vorzugeben und Ist-Positionen auszulesen. Hübel~\cite{huebel2022}
integrierte für eine automatisierte Ausschankvorrichtung zusätzliche externe
Sensorik (Füllstandsmessung) in die Zelle.

Beide Vorarbeiten haben gemeinsam, dass die Kommunikation zwischen PC und
Steuerung nicht echtzeitfähig im eigentlichen Sinn ist. 
Der MATLAB/KVP-Pfad läuft über eine zyklisch pollende TCP/IP-Verbindung unter Windows ohne garantierte Zykluszeiten, wie die im Rahmen dieses Projekts durchgeführte Baseline-Messung bestätigt (siehe Abschnitt~\ref{sec:baseline}). 
Für die im Thema~2 geforderte kraftgeregelte, reaktive Anwendung reicht dieser
Kommunikationspfad nicht aus.
Eine Kraftregelung mit kurzer Totzeit setzt
deterministische, garantierte Zykluszeiten voraus.

Ziel dieses Projekts ist es daher, die bestehende MATLAB/KVP-Anbindung auf
eine QUARC-basierte Simulink-Struktur mit deterministischem Abtastverhalten
zu portieren, einen Kraftsensor hard- und softwareseitig zu integrieren, die
End-to-End-Systemlatenzen des neuen Aufbaus zu erfassen und darauf
aufbauend eine totzeitkompensierte Regelungsapplikation zu realisieren. 
Der vorliegende Bericht beschreibt in Abschnitt~\ref{sec:aufgabenstellung} die
konkrete Aufgabenstellung, in Abschnitt~\ref{sec:sdt} den Stand der Technik
auf Basis der beiden Vorarbeiten sowie der Zellkomponenten, in
Abschnitt~\ref{sec:implementierung} die Umsetzung inklusive der
MATLAB/KVP-Baseline-Messung und in Abschnitt~\ref{sec:fazit} eine
zusammenfassende Bewertung und den weiteren Ausblick.
